\section{ReAct: intercalando Reasoning y Acting}

\subsection{Por qué no basta con solo pensar o solo actuar}

Chain-of-thought puede descomponer un problema, pero el razonamiento interno no puede corregir automáticamente conocimiento desactualizado o alucinaciones; la sola invocación de herramientas puede obtener información nueva, pero carece de planeación de largo plazo e integración entre observaciones. ReAct hace que el modelo genere de forma alternada:

\begin{itemize}
    \item \textbf{Thought}: el juicio actual, los subobjetivos y la justificación del siguiente paso;
    \item \textbf{Action}: una búsqueda, un clic, una consulta o una operación sobre el entorno;
    \item \textbf{Observation}: el estado externo devuelto por la herramienta.
\end{itemize}

\videofigure{figures/react-contribution.jpg}{ReAct intercala razonamiento en lenguaje y acciones sobre el entorno en una trayectoria unificada.}{00:04:40--00:06:46}

\videofigure{figures/react-loop.jpg}{Reasoning se encarga de planear e integrar, Action aporta evidencia nueva y corrige el razonamiento.}{00:06:46--00:07:32}

\subsection{Formalización como proceso de decisión secuencial}

En el paso de tiempo $t$, el agente elige una acción con base en el contexto histórico $c_t$:

$$
a_t\sim\pi_\theta(a\mid c_t),\qquad
c_{t+1}=c_t\oplus a_t\oplus o_{t+1},
$$

donde $o_{t+1}$ es la observation que devuelve el entorno después de ejecutar la acción. Lo particular de ReAct es que $a_t$ puede incluir tanto un thought en lenguaje como una action ejecutable, y ambos comparten el mismo contexto.

\videofigure{figures/react-setup.jpg}{El agente observa el entorno, elige una acción y escribe la nueva observation de vuelta en el contexto.}{00:07:32--00:09:06}

\begin{knowledgebox}{Thought también es una forma de comprimir el estado interno}
El modelo no necesita conservar todo el historial tal cual en un estado de programa estructurado; en cambio, resume en lenguaje natural las hipótesis, evidencias y pendientes actuales. Esto es flexible, pero también puede producir deriva de estado: un resumen erróneo temprano puede tomarse como hecho en el contexto posterior.
\end{knowledgebox}

\subsection{Reason-only, Act-only y ReAct}

\videofigure{figures/react-example.jpg}{En preguntas y respuestas de varios saltos, ReAct corrige el razonamiento interno mediante observations de búsqueda.}{00:09:06--00:17:18}

\begin{itemize}
    \item \textbf{Reason-only}: puede formar un plan, pero puede seguir razonando con base en una memoria errónea;
    \item \textbf{Act-only}: puede obtener hechos, pero tiende a carecer de gestión de subobjetivos e integración entre pasos;
    \item \textbf{ReAct}: usa el thought para decidir qué buscar, y la observation para verificar y actualizar el thought.
\end{itemize}

\subsection{Tareas intensivas en conocimiento: HotpotQA y FEVER}

HotpotQA exige responder preguntas de varios saltos a través de múltiples páginas de Wikipedia; FEVER exige recuperar evidencia para determinar si una afirmación es verdadera o falsa. Estas tareas evalúan no solo la respuesta final, sino también si se logra encontrar la evidencia que la sustenta.

\videofigure{figures/knowledge-tasks.jpg}{ReAct combina recuperación y razonamiento en tareas de preguntas y respuestas de varios saltos y de verificación de hechos.}{00:17:18--00:20:46}

La ventaja de ReAct es que la trayectoria es interpretable: los investigadores pueden ver qué buscó el modelo, en qué paso aceptó evidencia errónea y si repitió búsquedas. Pero un «thought visible» no garantiza que refleje fielmente el proceso causal interno del modelo; es más apropiado tratarlo como memoria de trabajo depurable.

\subsection{Tareas de decisión: WebShop}

WebShop exige que el agente navegue, filtre y compre productos en un sitio de compras simulado siguiendo las restricciones del usuario. El éxito no depende solo del conocimiento, sino también de operaciones con estado y de flujos de página que no pueden saltarse.

\videofigure{figures/webshop.jpg}{WebShop convierte restricciones en lenguaje natural en acciones web de varios pasos.}{00:20:46--00:22:10}

\subsection{Ventajas y modos de fallo}

\videofigure{figures/react-strengths.jpg}{La ventaja de ReAct es que es simple, interpretable y no requiere actualizar los parámetros del modelo.}{00:22:10--00:23:22}

Los fallos comunes incluyen:

\begin{itemize}
    \item la consulta de búsqueda es demasiado estrecha, y la primera evidencia errónea fija el rumbo de los pasos siguientes;
    \item la observation es demasiado larga, y la información clave queda sepultada por el ruido;
    \item el thought describe el estado de forma repetitiva sin hacer avanzar el plan;
    \item un error de herramienta o un cambio en la página invalida la trayectoria;
    \item falta una condición de detención, y el agente entra en bucle sobre la misma acción.
\end{itemize}

\begin{warningbox}{El resultado de una herramienta no es ground truth confiable por naturaleza}
Los resultados de búsqueda pueden estar desactualizados, una página web puede contener prompt injection malicioso y una base de datos puede tener permisos insuficientes. El agente debe registrar la fuente, verificar la evidencia y limitar los permisos de las herramientas, en lugar de elevar la observation a instrucción del sistema sin condiciones.
\end{warningbox}

\subsection{Resumen de la sección}

ReAct convierte el razonamiento en lenguaje de un proceso de generación cerrado en una trayectoria que puede interactuar con el entorno. Su principal aporte es la capacidad de autocorrección dentro de la tarea actual, pero los parámetros no se actualizan; para que el modelo mejore de forma duradera a partir de una gran cantidad de experiencias de ejecución, se necesita retroalimentación durante el entrenamiento.
